\documentclass[10pt,a4paper]{amsart}
\usepackage[margin=19mm]{geometry}
\usepackage{amsmath,amssymb,mathtools}
\usepackage[scr=rsfs,cal=euler]{mathalfa}
\usepackage{xfrac}
\usepackage[unicode,hidelinks,bookmarks=false]{hyperref}
\newcommand{\ie}{i.e.\ }
\newcommand{\cf}{cf.\ }
\newcommand{\resp}{respectively\ }
\newcommand{\Subsection}{\subsection}
\providecommand{\nobreakdash}{\nobreak\hbox{-}}
\setlength{\emergencystretch}{2em}
\multlinegap0pt
\usepackage{bm}
\usepackage{amssymb}
\usepackage{mathtools}
\usepackage{braket}
\usepackage{enumerate}
\usepackage[shortlabels]{enumitem}
\usepackage{csquotes}
\newtheorem{lem}{Lemma}
\newcommand{\cyclicpermutations}{\mathscr{C}(1,\dots, k, 2k, \dots, k+1)}
\title{NPA Hierarchy for Quantum Isomorphism and\texorpdfstring{\newline}{ } Homomorphism Indistinguishability$^{*}$}
\author{Prem Nigam Kar, David E. Roberson, Tim Seppelt, Peter Zeman}
\date{Source: arXiv:2407.10635v3 (CC BY 4.0)}
\begin{document}
\maketitle

\noindent\emph{Source and licence.} NPA Hierarchy for Quantum Isomorphism and\texorpdfstring{\newline}{ } Homomorphism Indistinguishability$^{*}$. Version arXiv:2407.10635v3 (CC BY 4.0), distributed by arXiv under the Creative Commons Attribution 4.0 International licence, http://creativecommons.org/licenses/by/4.0/, as recorded in the arXiv licence metadata for this exact version and shown on the abstract page https://arxiv.org/abs/2407.10635v3. Full licence notice: \texttt{licenses/arxiv-2407.10635-CC-BY-4.0.txt}.\par\smallskip

\noindent\textit{Editorial scope.} Counted unit: the article's lem lem:chain-collapse (source0.tex lines 1230--1233). The article's complete original proof of it is delivered with it (source0.tex lines 1234--1278, one \texttt{proof} environment of 45 lines). No mathematics is written, repaired or re-derived: the statement and the proof are the article's own lines, in the article's own notation, and no retained line is substituted or re-flowed.  No further source-local context is needed: every cross-reference of every retained line resolves inside the two delivered spans.  Nothing was omitted from any retained span, so the package carries no omission record.  No retained line contains a citation, so the wrapper carries no bibliography and no bibliography entry.  No retained line contains a figure environment, an image-inclusion command or drawing code, so no image is delivered and the package declares no asset dependency.  Editorial formatting only, in addition: an AMS \texttt{amsart} preamble that restates the article's own declarations and macros used by the retained lines, namely the environments \texttt{lem} on separate counters, together with the standard packages those lines need.  The retained environments print in this document as Lemma 1 in order of appearance, whereas the article prints the same items with the numbering of its own full document.  The complete native source of the article as distributed in the arXiv e-print for this exact version is bundled unchanged as \texttt{sources/162334/source0.tex}.
\begin{lem} \label{lem:chain-collapse}
	If $l \geq 1$ is such that $S_l = S_{l+1}$,
	then $S_l = S$.
\end{lem}

\begin{proof}
	Since $S_l \subseteq S$, it suffices to show the converse inclusion.
	We show by structural induction on $\boldsymbol{F} \in \mathcal{P}_k$ that $\boldsymbol{F}_G \oplus \boldsymbol{F}_H \in S_l$.
	If $\boldsymbol{F}$ is atomic then the claim follows immediately.
	For the inductive step, distinguish cases:
	\begin{enumerate}
		\item There exist $\boldsymbol{F}' \in \mathcal{P}_k$ and $\boldsymbol{Q} \in \mathcal{Q}_k^P$ such that $\boldsymbol{F} = \boldsymbol{Q} \odot \boldsymbol{F}'$.
		
		By the inductive hypothesis, $\boldsymbol{F}'_G \oplus \boldsymbol{F}'_H \in S_l$.
		Hence, there exist $\boldsymbol{F}^i \in \mathcal{P}_k$ with $\nu(\boldsymbol{F}^i) \leq l$ and coefficients $\alpha_i \in \mathbb{R}$ such that $\boldsymbol{F}'_G \oplus \boldsymbol{F}'_H  = \sum \alpha_i \boldsymbol{F}^i_G \oplus \boldsymbol{F}^i_H$.
		By definition, $\nu(\boldsymbol{Q} \odot \boldsymbol{F}^i) \leq \nu(\boldsymbol{F}^i) \leq l$.
		Hence,
		\[
			\boldsymbol{F}_G \oplus \boldsymbol{F}_H = (\boldsymbol{Q}_G \oplus \boldsymbol{Q}_H) \odot \sum \alpha_i \boldsymbol{F}^i_G \oplus \boldsymbol{F}^i_H
			= \sum \alpha_i (\boldsymbol{Q} \odot \boldsymbol{F}^i)_G
			 \oplus (\boldsymbol{Q} \odot \boldsymbol{F}^i)_H
			 \in S_l.
		\]
		
		\item There exist $\boldsymbol{F}', \boldsymbol{F}'' \in \mathcal{P}_k$ such that $\boldsymbol{F} = \boldsymbol{F}' \cdot \boldsymbol{F}''$.
		
		By the inductive hypothesis, $\boldsymbol{F}'_G \oplus \boldsymbol{F}'_H,\boldsymbol{F}''_G \oplus \boldsymbol{F}''_H  \in S_l$.
		Hence, there exist $\boldsymbol{F}^i \in \mathcal{P}_k$, $\boldsymbol{K}^j \in \mathcal{P}_k$ with $\nu(\boldsymbol{F}^i),\nu(\boldsymbol{K}^j) \leq l$ and coefficients $\alpha_i, \beta_j \in \mathbb{R}$ such that $\boldsymbol{F}'_G \oplus \boldsymbol{F}'_H  = \sum \alpha_i \boldsymbol{F}^i_G \oplus \boldsymbol{F}^i_H$
		and $\boldsymbol{F}''_G \oplus \boldsymbol{F}''_H  = \sum \beta_j \boldsymbol{K}^j_G \oplus \boldsymbol{K}^j_H$.
		By definition, $\nu(\boldsymbol{F}^i \cdot \boldsymbol{K}^j) \leq \max\{ \nu(\boldsymbol{F}^i), \nu(\boldsymbol{K}^j)\} \leq l+1$ for all $i, j$.
		Hence, 
		\[
			\boldsymbol{F}_G \oplus \boldsymbol{F}_H = \sum \alpha_i \beta_j (\boldsymbol{F}^i \cdot \boldsymbol{K}^j)_G \oplus 
			(\boldsymbol{F}^i \cdot \boldsymbol{K}^j)_H
			\in S_{l+1}  = S_l.
		\]
		
		\item There exist $\boldsymbol{F}' \in \mathcal{P}_k$ and $\sigma \in \cyclicpermutations$ such that $\boldsymbol{F} = (\boldsymbol{F}')^\sigma$.
		
		By the inductive hypothesis, $\boldsymbol{F}'_G \oplus \boldsymbol{F}'_H \in S_l$.
		Hence, there exist $\boldsymbol{F}^i \in \mathcal{P}_k$ with $\nu(\boldsymbol{F}^i) \leq l$ and coefficients $\alpha_i \in \mathbb{R}$ such that $\boldsymbol{F}'_G \oplus \boldsymbol{F}'_H  = \sum \alpha_i \boldsymbol{F}^i_G \oplus \boldsymbol{F}^i_H$.
		By definition, $\nu((\boldsymbol{F}^i)^\sigma) \leq \nu(\boldsymbol{F}^i) \leq l$.
		Hence,
		\[
		\boldsymbol{F}_G \oplus \boldsymbol{F}_H = \left( \sum \alpha_i \boldsymbol{F}^i_G \oplus \boldsymbol{F}^i_H \right)^\sigma
		= \sum \alpha_i ((\boldsymbol{F}^i)^\sigma)_G \oplus ((\boldsymbol{F}^i)^\sigma)_H
		\in S_l. \qedhere
		\] 
	\end{enumerate} 
\end{proof}

\end{document}
